% Fuentes, tamaño, diccionario mínimo de las 10 variables principales, muestra o tablas de nulos, outliers, problemas de calidad y decisiones de limpieza ya tomadas o pendientes.

\section{Relevamiento de Fuentes de Datos a utilizar}

El conjunto de datos utilizado en este proyecto fue obtenido de Kaggle y corresponde al dataset \textit{BankSim}. Este conjunto fue desarrollado por \textcite{inproceedings} con el objetivo de simular, de manera sintética, un sistema de pagos bancarios que reproduzca características propias de transacciones financieras reales.

El dataset contiene información sobre las transacciones realizadas durante la simulación y se encuentra compuesto por diez variables, cuya descripción y tipo conceptual se presentan en el Cuadro~\ref{tab:variables_dataset}.

\begin{table}[!htbp]
    \centering
    \caption{Variables incluidas en el dataset BankSim}
    \label{tab:variables_dataset}
    \begin{tabular}{lll}
        \toprule
        \textit{Variable} & \textit{Descripción} & \textit{Tipo conceptual} \\
        \midrule
        \texttt{step}        & Día de la simulación                    & Temporal \\
        \texttt{customer}    & Identificador del cliente               & Categórica / ID \\
        \texttt{age}         & Rango de edad del cliente               & Categórica ordinal \\
        \texttt{gender}      & Género o tipo del cliente               & Categórica \\
        \texttt{zipcodeOri}  & Código postal de origen del cliente     & Categórica \\
        \texttt{merchant}    & Identificador del comercio              & Categórica / ID \\
        \texttt{zipMerchant} & Código postal del comercio              & Categórica \\
        \texttt{category}    & Categoría del comercio o de la compra   & Categórica \\
        \texttt{amount}      & Importe de la transacción               & Numérica \\
        \texttt{fraud}       & Indica si la transacción es fraudulenta & Variable objetivo binaria \\
        \bottomrule
    \end{tabular}
\end{table}

\begin{nota}
El dataset no contiene nulos ni outliers debido a errores en la carga de datos.
\end{nota}

\subsection{Tratamiento de los datos}

\paragraph{Columnas eliminadas}

Las columnas \texttt{age} y \texttt{gender} fueron eliminadas debido a que los autores de \textit{BankSim} indican que estas variables no fueron utilizadas para modelar el comportamiento de consumo de los agentes, los autores señalan: ``However, we did not programme the consumption behaviour of agents based on gender and age'' \parencite[p. 150]{inproceedings}.

Asimismo, las columnas \texttt{zipcodeOri} y \texttt{zipMerchant} fueron eliminadas porque contienen un único valor en todo el dataset y, por lo tanto, no aportan información al modelo.

\paragraph{Columnas transformadas}

Mediante el \textit{feature engineering} propuesto por \textcite[sección 3.3]{leborgne2022fraud}, se generaron las variables \nolinkurl{tx_count_last_}\(n\)\nolinkurl{d} y \nolinkurl{avg_amount_last_}\(n\)\nolinkurl{d} para representar la actividad reciente de cada cliente. Asimismo, se crearon \nolinkurl{merchant_tx_delayed_}\(n\)\nolinkurl{d}, \nolinkurl{merchant_risk_}\(n\)\nolinkurl{d}, \nolinkurl{category_tx_delayed_}\(n\)\nolinkurl{d} y \nolinkurl{category_risk_}\(n\)\nolinkurl{d}, considerando un retraso de siete días en la disponibilidad de la información sobre fraude. En todos los casos, \(n \in \{1,7,30\}\) días.

\begin{nota}
Luego de las transformaciones, el dataset consta de 24 columnas, sin haber aplicado aún \textit{one-hot encoding}.
\end{nota}